% ============================================================================
%  SECCIÓN 3.23: DISEÑO RESPONSIVE
% ============================================================================

\section{Dise\~no responsive}

El dise\~no responsive garantiza que la interfaz de la aplicaci\'on se adapte
correctamente a diferentes tama\~nos de pantalla y condiciones de uso
\citep{nielsen1994, pressman2010}. En el contexto del desarrollo m\'ovil, esto
implica considerar la aparici\'on del teclado virtual, la variaci\'on en los
tama\~nos de dispositivo y la orientaci\'on de la pantalla.

\subsection{Adaptaci\'on al teclado virtual}

Las pantallas que contienen formularios (login y registro) utilizan el
componente \texttt{KeyboardAvoidingView} de React Native, que detecta
autom\'aticamente la aparici\'on del teclado virtual y ajusta el layout de la
interfaz para que los campos de texto no queden ocultos detr\'as del teclado.

\begin{lstlisting}
<KeyboardAvoidingView
  style={styles.keyboardView}
  behavior={Platform.OS === 'ios' ? 'padding' : 'height'}
>
  <ScrollView
    contentContainerStyle={styles.scrollContent}
    keyboardShouldPersistTaps="handled"
  >
    {/* Contenido del formulario */}
  </ScrollView>
</KeyboardAvoidingView>
\end{lstlisting}

La propiedad \texttt{behavior} se configura de forma condicional seg\'un la
plataforma: en iOS se utiliza \texttt{'padding'} y en Android se utiliza
\texttt{'height'}. Adem\'as, \texttt{keyboardShouldPersistTaps="handled"}
permite que los toques en botones se registren aunque el teclado est\'e
visible.

\subsection{Desplazamiento con ScrollView}

El componente \texttt{ScrollView} se utiliza en las pantallas de login y
registro para permitir el desplazamiento del contenido cuando el formulario
excede el tama\~no de la pantalla. Esto es especialmente importante en
dispositivos con pantallas peque\~nas o cuando el teclado virtual reduce el
espacio visible.

\begin{lstlisting}
<ScrollView
  contentContainerStyle={styles.scrollContent}
  keyboardShouldPersistTaps="handled"
>
  <View style={styles.container}>
    {/* Contenido del formulario */}
  </View>
</ScrollView>
\end{lstlisting}

\subsection{Layouts flexibles con Flexbox}

El sistema Flexbox de React Native se utiliza de forma extensiva para crear
layouts que se adaptan autom\'aticamente al tama\~no de la pantalla. Las
propiedades clave utilizadas son:

\begin{table}[H]
  \centering
  \caption{Propiedades Flexbox utilizadas en el dise\~no responsive}
  \contenidoConFuente{%
    \begin{tablaAPA}
      \begin{tabular}{@{}p{3.4cm}p{12.2cm}@{}}
        \toprule
        \textbf{Propiedad} & \textbf{Aplicaci\'on} \\\
        \midrule
        \texttt{flex: 1} & Permite que un contenedor ocupe todo el espacio
        disponible, distribuyendo el espacio proporcionalmente entre los
        elementos hijos. \\\
        \addlinespace
        \texttt{flexDirection: 'column'} & Organiza los elementos de arriba
        hacia abajo (layout vertical), patr\'on principal de las pantallas. \\\
        \addlinespace
        \texttt{flexDirection: 'row'} & Organiza los elementos de izquierda a
        derecha (layout horizontal), utilizado en footers y barras de
        herramientas. \\\
        \addlinespace
        \texttt{alignItems: 'center'} & Centra los elementos hijos
        horizontalmente dentro de su contenedor. \\\
        \addlinespace
        \texttt{justifyContent: 'center'} & Centra los elementos hijos
        verticalmente dentro de su contenedor. \\\
        \addlinespace
        \texttt{justifyContent: 'space-between'} & Distribuye el espacio
        uniformemente entre los elementos, utilizado en el splash. \\\
        \bottomrule
      \end{tabular}
    \end{tablaAPA}%
  }{Elaboraci\'on propia en base al c\'odigo fuente de la aplicaci\'on.}
\end{table}

\subsection{Uso de porcentajes y anchos relativos}

La aplicaci\'on utiliza anchos relativos (porcentajes) para componentes que
deben adaptarse al tama\~no de la pantalla, como el bot\'on \textit{``Comenzar''}
del splash que ocupa el 100\% del ancho disponible:

\begin{lstlisting}
footer: {
  width: '100%',
  zIndex: 1,
},
\end{lstlisting}

\subsection{Pr\'acticas de accesibilidad visual}

Se aplicaron las siguientes pr\'acticas para mejorar la accesibilidad visual de
la interfaz:

\begin{itemize}
  \item Contraste adecuado entre texto y fondo, siguiendo las pautas de
    accesibilidad (texto principal \texttt{\#0F172A} sobre fondo
    \texttt{\#F8FAFC}, ratio de contraste superior a 4.5:1).
  \item Tama\~nos de texto m\'inimos de 12\,pt para garantizar la legibilidad.
  \item Campos de formulario con altura m\'inima de 44\,pt para facilitar la
    interacci\'on con el dedo.
  \item Uso de \texttt{placeholderTextColor} para diferenciar el texto
    placeholder del texto ingresado por el usuario.
  \item Indicadores visuales de estado: bordes rojos para errores de
    validaci\'on, indicadores de carga durante el env\'io de formularios.
\end{itemize}
